\documentclass[conference]{IEEEtran}
\IEEEoverridecommandlockouts
\usepackage{cite}
\usepackage{amsmath,amssymb,amsfonts}
\usepackage{graphicx}
\usepackage{booktabs}
\usepackage{array}
\newcolumntype{L}[1]{>{\raggedright\arraybackslash}p{#1}}
\usepackage{algorithm}
\usepackage{algpseudocode}
\usepackage{textcomp}
\usepackage{xcolor}
\usepackage{tikz}
\usetikzlibrary{positioning,arrows.meta,fit}
\usepackage{url}
\newcommand{\todo}[1]{\textcolor{red}{[TODO: #1]}}
\newcommand{\kbgym}{KBGym}

\begin{document}

\title{Training a Knowledge Base:\\Supervised Structure Learning for Agent-Curated Document Stores\\
\thanks{Draft v0.2 --- August 2026. Results marked preliminary are being re-established on the final environment.}}

\author{\IEEEauthorblockN{Anonymous Author}
\IEEEauthorblockA{\textit{Affiliation withheld} \\
email withheld}}

\maketitle

\begin{abstract}
Retrieval-augmented generation treats the document store as a frozen input. A recent
wave of systems instead lets an LLM agent continuously curate a persistent knowledge
base, yet none of this work measures what curation does to the store itself. We invert
the usual framing: \emph{the knowledge base is the model}. A training agent consumes
supervised QA pairs one at a time---attempting each question against the current store,
receiving the gold answer, then consolidating verified reasoning into the store's
structure. An independent frozen-store reader is then
examined on held-out questions under a fixed action budget. Our central hypothesis is an \emph{efficiency} claim:
by materializing reasoning paths into index and shortcut structure, a trained store
lets an unchanged reader answer at \emph{non-inferior accuracy in fewer steps},
amortizing one-time training cost across all future queries. The protocol yields the
full apparatus of machine learning for knowledge bases: learning curves, a train/test
generalization gap, per-class transfer, and a break-even analysis of training cost
against per-question savings. On a contamination-free synthetic universe with
document-level ground truth (500 people, 5{,}864 atomic documents), two
training epochs lift \emph{re-encountered} questions from 0.59 to 0.85 F1
at 38\% fewer steps, with gains concentrated exactly in the relation-level
question classes the design predicts, and the links the agent builds join
documents that co-occur in a question's support set at 720$\times$ the
chance rate. On \emph{unseen} questions, however, the efficiency claim
fails: held-out accuracy is non-inferior to a flat store at two action
budgets but never cheaper ($\rho \approx 1$), and neither do
offline-construction baselines (GraphRAG- and HippoRAG2-style) save any
steps. A difficulty-stratified attribution locates the reason: built
structure pays where derivation is expensive and taxes retrieval where it
is not. We report the negative result with the apparatus that makes it
legible.
\end{abstract}

\begin{IEEEkeywords}
knowledge bases, LLM agents, retrieval-augmented generation, agentic RAG, memory, benchmarks
\end{IEEEkeywords}

\section{Introduction}
The dominant way to give language models durable knowledge is retrieval-augmented
generation (RAG): embed a corpus once, retrieve at query time, never write back. A
rapidly growing line of work breaks this asymmetry. Agent-maintained wikis compile
sources into evolving page networks~\cite{llmwiki-rar,llmwiki-stream,karpathy-llmwiki};
agentic memory systems accumulate and link notes across sessions~\cite{amem}; and
``sleep-time'' agents reorganize memory between conversations~\cite{letta-sleep}. The
pattern has reached practice ahead of measurement: the community's own description of
the wiki pattern lists a \emph{lint} step---checking contradictions, stale claims,
orphan pages, missing cross-links---as an engineering habit with no formal evaluation
attached~\cite{karpathy-llmwiki}, and the closest published systems evaluate only
downstream answer accuracy, never the store~\cite{llmwiki-rar,llmwiki-stream}.

We study the store directly, under the oldest framing machine learning has:
\emph{train/test}. The knowledge base is the model. Supervised questions arrive one at
a time; a training agent answers each from the current store, is graded, and then, in
an explicit consolidation phase, edits the store's structure so that the question's
\emph{class} becomes easier for a future reader. Evaluation freezes the store and
examines an independent reader with no gold access and a fixed action budget on
held-out questions.

The value produced by such training is best understood as \emph{index construction}:
a database does not change its data when an index is built, yet queries get cheaper.
Each training question's verified reasoning path is materialized as
\emph{access structure}---a multi-hop chain becomes a traversable link
path, a key that search handles badly gains an index document linking
everything under it---so future readers pay a lookup price instead of a
re-derivation price. The distinction from writing the answer into a
document is the whole point: an answer serves the instance that produced
it, an index serves every question that enters through the same key. Training cost is thereby amortized: after a computable number of questions, the
one-time cost of curation is repaid by per-question savings.

\textbf{Contributions.}
\begin{enumerate}
\item \textbf{Agent-trained knowledge bases.} We propose letting an LLM agent
\emph{train} a document store under the standard machine-learning training
paradigm, with the knowledge base playing the role of the model: supervised
(question, answer) pairs are consumed one at a time over multiple epochs; each
iteration runs a forward pass (answer the question from the current store) and
a backward pass (with the revealed gold in hand, \emph{edit the store
itself}---adding, rewriting, deleting, and linking documents; building index
and navigation documents---so the question's class becomes easier for a future
reader); held-out performance is the loss. The knowledge base, not the model
weights, is the learned object, and the structure is carved incrementally by
the questions actually asked rather than induced from the corpus alone. In
machine-learning terms the contrast with offline graph construction
(GraphRAG, RAPTOR, HippoRAG) is exact: those are \emph{unsupervised} ---
structure inferred from the corpus with no labels and no notion of what will
be asked --- while ours is \emph{supervised}: (question, answer) pairs are
the labels, and structure is optimized against verified answers.
\item \textbf{A train/test protocol for knowledge bases}: two-phase
(forward/backward) training iterations, single-shot gold reveal,
frozen-store examination under action budgets, learning curves, and a
three-layer generalization decomposition (re-encountered instances / unseen
instances / unseen templates).
\item \textbf{\kbgym}, a contamination-free environment in the PhantomWiki
style~\cite{phantomwiki}: a fictional-person universe rendered as a graph of
atomic single-sentence documents; ten diagnostic question classes with exact
answers and support sets; deterministic, integer-exact store metrics; no LLM
judge anywhere.
\item \textbf{Cost as a first-class outcome}: accuracy, store size,
training cost, and per-question answering cost reported together, with a
break-even analysis of when training a store pays for itself.
\end{enumerate}

\section{Related Work}
\subsection{Retrieval-Augmented Generation}
Classic RAG retrieves from a frozen corpus with a learned dense retriever and
generates conditioned on the results~\cite{rag-lewis,dpr,realm,fid}. The
corpus is an input, never an output: nothing the system learns at answer time
flows back into the store. Our work targets exactly this missing backward
edge.

\subsection{Agent-Maintained Knowledge Bases}
Self-evolving LLM-wiki retrieval~\cite{llmwiki-rar} has a single agent write and
update wiki pages during question answering, outperforming static RAG on multi-hop
benchmarks; streaming knowledge compilation~\cite{llmwiki-stream} maintains a
time-evolving wiki under a document stream. STORM~\cite{storm} generates
Wikipedia-style articles with multi-perspective research agents---one-shot
authorship rather than continual maintenance. The published systems evaluate
downstream task scores only:
the wiki itself is never measured, there is no train/test split, and the evaluation
corpora (HotpotQA, MuSiQue) are parametrically contaminated. Our protocol adds the
missing measurement layer.

\subsection{Agentic Retrieval}
Agentic RAG moves retrieval into the reasoning loop: ReAct interleaves
reasoning with tool calls~\cite{react}, IRCoT interleaves chain-of-thought
with retrieval~\cite{ircot}, Self-Ask decomposes into sub-questions before
searching~\cite{selfask}, FLARE retrieves when generation confidence
drops~\cite{flare}, Self-RAG learns when to retrieve and how to judge what
returns~\cite{selfrag}, and recent systems frame the retrieve-or-reason
decision as an MDP (DeepRAG~\cite{deeprag}) or train the search policy with
reinforcement learning (Search-R1~\cite{searchr1}; surveyed
in~\cite{agentic-rag-survey}). Our reader \emph{is} an agentic-RAG loop---budgeted
iterative search-and-read. The distinction is which side of the loop learns: that line
improves the \emph{query-side policy} against a fixed corpus; we hold the policy fixed
and train the \emph{store}. The axes are orthogonal and composable; our budget sweep
measures exactly what a better store buys an unchanged agentic reader.

\subsection{Offline Structure Construction}
GraphRAG~\cite{graphrag}, RAPTOR~\cite{raptor}, LightRAG~\cite{lightrag}, and
HippoRAG~\cite{hipporag,hipporag2} build entity graphs, summary hierarchies,
or PageRank-linked knowledge graphs over a corpus \emph{before} any question
arrives. These are our strongest baselines and the natural
contrast: their structure is unsupervised and built blind to the query distribution;
ours is carved incrementally by the questions actually asked.

\subsection{Agent Memory and Experience Accumulation}
A rich line lets agents accumulate experience across episodes: MemGPT pages
context in and out of external memory~\cite{memgpt}; Mem0 and A-MEM maintain
linked long-term note stores~\cite{mem0,amem}; generative agents retrieve and
reflect over an event stream~\cite{generative-agents}; Voyager grows a skill
library~\cite{voyager}; Reflexion, ExpeL, and Agent Workflow Memory distill
verbal feedback, cross-task insights, and reusable workflows,
respectively~\cite{reflexion,expel,awm}; Dynamic Cheatsheet and ACE evolve a
test-time strategy memory or context playbook~\cite{cheatsheet,ace}. In all
of these the artifact is a prompt or a private memory serving one agent; ours
is a shared, inspectable document store whose content is
ground-truth-checkable and whose curation quality is itself the measured
outcome.

\subsection{Memory Evaluation and Contamination Control}
LongMemEval and successors probe assistant memory with QA over past sessions,
typically via LLM judges~\cite{longmemeval}; LoCoMo evaluates very-long-term
conversational memory~\cite{locomo}; MemDelta documents hidden confounds and
argues for controlled baselines~\cite{memdelta}, supporting our no-store
baseline discipline. Multi-hop QA datasets---HotpotQA~\cite{hotpotqa},
MuSiQue~\cite{musique}, 2WikiMultiHopQA~\cite{twowiki}---supply real-text
questions with annotated support, but their public corpora are parametrically
contaminated for current models. PhantomWiki~\cite{phantomwiki} generates
contamination-free fictional universes with programmatically exact answers;
our generator follows its recipe, and SQuAD-style normalization~\cite{squad}
provides the grading.

\subsection{Knowledge Editing}
ROME and MEMIT edit factual associations \emph{inside model
weights}~\cite{rome,memit}. We pursue the non-parametric complement: facts
live in an external store the agent can inspect, reorganize, and be audited
on---no gradient access, full provenance, and edits that transfer across
models. STaR-style bootstrapping~\cite{star} keeps only verified reasoning
for further training; our forward/backward iteration inherits that
predict-then-learn shape, with the store rather than the weights as the
learned object.

\definecolor{okblue}{HTML}{0072B2}
\definecolor{okorange}{HTML}{E69F00}
\definecolor{okverm}{HTML}{D55E00}
\definecolor{okgreen}{HTML}{009E73}
\definecolor{okpurple}{HTML}{9A5EA8}
\definecolor{okteal}{HTML}{56B4E9}

\begin{figure*}[t]
\centering
\begin{tikzpicture}[
  font=\small,
  node/.style={circle, draw=okblue!80, fill=okblue!25, inner sep=1.6pt},
  nav/.style={circle, draw=okverm, fill=okverm!85, inner sep=2.2pt},
  stage/.style={draw, rounded corners=2pt, align=center, inner sep=5pt},
  lbl/.style={font=\footnotesize\itshape, text=black!70},
  arr/.style={-{Stealth[length=2.2mm]}, thick, draw=black!65},
  >=Stealth]

% ---- left: initial store
\node[stage, draw=okblue!70, fill=okblue!5, minimum width=3.4cm,
      minimum height=2.5cm] (K0) {};
\node[above=1mm of K0.north, anchor=south] {\textbf{initial store} $K_0$};
\node[lbl, below=0mm of K0.south, anchor=north] {atomic documents, no links};
\foreach \i/\x/\y in {1/-1.1/0.6, 2/-0.3/0.8, 3/0.6/0.55, 4/-1.05/-0.1,
                      5/0.1/0.05, 6/1.05/0.0, 7/-0.6/-0.7, 8/0.35/-0.75, 9/1.0/-0.65}
  \node[node] (a\i) at ([shift={(\x,\y)}]K0.center) {};

% ---- middle: training loop
\node[stage, draw=okorange!85, fill=okorange!6, right=1.5cm of K0, minimum width=4.2cm, align=left] (train) {
  \textbf{training iteration} $(q, a^{*})$\\[1pt]
  \textcolor{okblue!80!black}{\emph{Phase 1} forward ($N_1$): search/read $\to$ \texttt{answer}}\\
  \hspace{2mm}$\hookrightarrow$ result reveals $a^{*}$, F1\\[1pt]
  \textcolor{okverm!90!black}{\emph{Phase 2} backward ($N_2$): add/edit/delete/}\\
  \textcolor{okverm!90!black}{\hspace{2mm}link/unlink $\to$ \texttt{done}}
};
\node[lbl, below=0mm of train.south, anchor=north]
  {$\times$ 150 questions $\times$ epochs};

% ---- right: trained store
\node[stage, draw=okgreen!80, fill=okgreen!5, right=1.5cm of train, minimum width=3.4cm, minimum height=2.5cm] (KT) {};
\node[above=1mm of KT.north, anchor=south] {\textbf{trained store} $K_T$};
\node[lbl, below=0mm of KT.south, anchor=north] {links + navigation documents};
\foreach \i/\x/\y in {1/-1.1/0.6, 2/-0.3/0.8, 3/0.6/0.55, 4/-1.05/-0.1,
                      5/0.1/0.05, 6/1.05/0.0, 7/-0.6/-0.7, 8/0.35/-0.75, 9/1.0/-0.65}
  \node[node] (b\i) at ([shift={(\x,\y)}]KT.center) {};
\node[nav] (hub) at ([shift={(-0.2,-0.35)}]KT.center) {};
\draw[okgreen!70!black, thick] (hub)--(b4) (hub)--(b5) (hub)--(b7) (hub)--(b8)
                (b2)--(b5) (b3)--(b6) (b6)--(b9);

% ---- arrows
\draw[arr] (K0) -- node[above, lbl]{} (train);
\draw[arr] (train) -- (KT);

% ---- bottom: exam
\node[stage, draw=okpurple!80, fill=okpurple!6, below=0.9cm of KT, minimum width=4.6cm, align=center] (exam)
  {\textbf{exam} (store frozen)\\ fresh reader: search/read/\texttt{answer},
   budget $M$\\ held-out questions, gold never shown};
\draw[arr] (KT) -- (exam);
\node[stage, draw=okpurple!80, fill=okpurple!6, left=1.2cm of exam, align=center] (out)
  {F1 (non-inferior?)\\ steps per question ($\rho<1$?)};
\draw[arr] (exam) -- (out);
\end{tikzpicture}
\caption{The train/test protocol. A training agent consumes supervised QA pairs,
answering each from the current store (gold revealed only after its single answer)
and consolidating verified reasoning into structure. Evaluation freezes the store and
examines an independent budgeted reader on held-out questions; the hypothesis
(Eq.~\ref{eq:hypothesis}) is non-inferior accuracy at strictly lower per-question
cost.}
\label{fig:protocol}
\end{figure*}

\section{Problem Formulation}
A knowledge base $K$ is a directed graph of \emph{documents}. A document is
deliberately small --- one self-contained sentence --- so that what the store
learns is visible in its structure rather than hidden inside prose, and so
that every edit is attributable. Concretely it carries six fields:

\begin{itemize}\itemsep2pt
\item \texttt{id} --- a stable handle; links and every log entry refer to it.
\item \texttt{text} --- one sentence stating one fact, with full names and no
pronouns, so it is self-contained and its embedding is about that fact alone.
It is the only thing \texttt{search} matches against and the only thing a
reader is shown.
\item \texttt{links} --- outgoing edges to other documents. \texttt{read}
returns a document together with the full text of everything it links, so a
document's edges are what makes traversal cheaper than search.
\item \texttt{origin} --- which universe fact this document represents, or
\texttt{None} when the agent wrote it. This is what makes the coverage and
duplication ledger exact: every original fact is either still represented or
demonstrably lost.
\item \texttt{flag} --- \texttt{None} for an untouched original,
\texttt{authored} for a document the agent wrote, \texttt{edited} for an
original it rewrote. Structure the agent added is therefore separable from
the corpus it was given, which is what the structural measurements report on.
\item \texttt{absorbed} --- origins folded in when a duplicate was merged, so
a merge loses no provenance.
\end{itemize}

An \emph{index document} is not a distinct type: it is a document whose text
names a key and whose value lies in its links rather than in the fact it
states. Nothing in the store distinguishes the two, which is why an index
that carries no links is indistinguishable from a sentence nobody needs. The \emph{reader} is a fixed
LLM policy with three tools---vector \texttt{search}, document \texttt{read} (text plus
linked documents), and \texttt{answer}---a fixed action budget $M$, and a short-term
memory holding only the current question and the last $k$ action--result pairs.
Reader performance on a question distribution $Q$ defines store quality along two
axes:
\begin{align*}
\mathrm{acc}(K) &= \mathbb{E}_{q\sim Q}\!\left[\mathrm{F1}\big(\mathrm{reader}(K,q)\big)\right],\\
\mathrm{cost}(K) &= \mathbb{E}_{q\sim Q}\!\left[\mathrm{steps}\big(\mathrm{reader}(K,q)\big)\right].
\end{align*}
\emph{Training} is any procedure that consumes supervised pairs $(q, a^{*})$ and edits
$K$, producing $K_T$ from the initial store $K_0$. What such a procedure is
meant to learn is the store's \emph{access structure}, and it is worth
being precise about this, because the action set admits a shortcut that
looks like learning and is not. A store can be changed in two ways: by
adding \emph{content} --- a document stating a fact, including a fact the
store already implies, such as the answer to a question just solved --- or
by adding \emph{access structure}: index documents whose value lies in the
links they carry to the documents belonging to some key, and edges between
documents that a chain must step across (including edges between index
documents, which give navigation its levels). Only the second generalizes.
A document recording a verified answer serves the one instance that
produced it, and the instance will not recur; an index over a key serves
every question whose first search names that key, and an edge serves every
chain that must make that hop, whatever entity it concerns. Our hypothesis
is therefore a claim about access structure specifically, and we measure
the store's navigability --- index out-degree and the fraction of content
reachable from an index in one read --- alongside accuracy and cost. Our central hypothesis is a joint
condition on the two axes:
\begin{equation}
\mathrm{acc}(K_T) \;\ge\; \mathrm{acc}(K_0) - \delta
\quad\text{and}\quad
\rho \;=\; \frac{\mathrm{cost}(K_T)}{\mathrm{cost}(K_0)} \;<\; 1,
\label{eq:hypothesis}
\end{equation}
with non-inferiority margin $\delta = 0.03$. Training cost amortizes: if training
spent $C_{\mathrm{train}}$ tokens and the per-question saving is
$c_0 - c_T$ tokens, the \emph{break-even point}
\begin{equation}
N^{*} \;=\; \frac{C_{\mathrm{train}}}{c_0 - c_T}
\label{eq:breakeven}
\end{equation}
is the query volume after which curation has paid for itself---the single number
that tells a practitioner whether training a store is worth it for their workload.

The correspondence to supervised learning is exact: the store configuration is the
parameter vector, one training iteration is one optimization step, the training agent
is the optimizer, and held-out reader performance is the loss. Formally, the
environment with a \emph{fixed} reader policy induces a Markov process over store
configurations; training is a guided walk through configuration space, and the
questions are its supervision signal. Because training-agent forward passes and test
readers use identical tools, budgets, and the single-shot answer rule, the train/test
generalization gap is well defined: the three-layer decomposition compares accuracy on
re-encountered train instances (retention carried \emph{only} by the store---the
agent's memory resets every question), unseen instances of trained templates
(class-level generalization), and reserved templates (within-category transfer).

\section{The \kbgym{} Environment}\label{sec:env}
\subsection{Universe and Store}
A generator in the PhantomWiki style samples a fictional population (500 people in
the main configuration): family trees, spouses, friendships, per-person attributes
(job, hobby, city), globally distinct birthdates, and \emph{name distractors}---extra
people sharing first names with questioned subjects. Facts render through fixed
templates into $\sim$5{,}900 atomic single-sentence documents (``Alice Johnson's job is
arborist.''). The initial store has \emph{zero links}: a bag of atomic facts.
Document embeddings back \texttt{search} and are maintained
automatically as documents are added, edited, or deleted.

\subsection{Actions}\label{sec:actions}
Reader and training-forward mode: \texttt{search(query)} (top-$k$ document ids
with full text), \texttt{read(id)} (text plus links, each with target text),
\texttt{answer(text)}. Training consolidation adds: \texttt{add(text)},
\texttt{edit(id,text)}, \texttt{delete(id)}, \texttt{link(a,b)},
\texttt{unlink(a,b)}, \texttt{done()}. Provenance is tracked throughout: initial
documents carry an origin; edited ones are flagged; authored ones have no origin.
Coverage and duplication are integer-exact; any authored content is attributable.

\subsection{Question Classes and Splits}
Ten categories over $\sim$26 templates. \emph{Forward} classes QC1--QC3 (1--3-hop
named-entity chains) are natively easy for document-level retrieval and serve as the
non-inferiority floor. \emph{Hard} classes target relations \emph{between}
documents: aggregation counts (QC4, QC8), abstention (QC5), multi-constraint joins
without a name anchor (QC6), set intersection (QC7), comparison/superlatives over
birthdates (QC9), and reverse lookup with a uniqueness guarantee (QC10). Golds are
computed by graph traversal over the fact layer; every question carries template id,
gold variants, and its support set. Splits (instance-disjoint): \texttt{train} (150;
i.i.d.\ over all categories, 60\% hard), \texttt{test\_in} (100; unseen instances of
trained templates), \texttt{test\_out} (50; one reserved template per category),
\texttt{eval} (30; stratified, drives the per-epoch learning curve so the final test
is touched exactly once). Grading is SQuAD-normalized F1/EM.

\subsection{Design Rationale}\label{sec:design}
Four decisions carry the methodology. \emph{(i) No LLM judge anywhere}: golds are
exact, grading is SQuAD normalization, store metrics are integer counts over
provenance---every number in this paper is reproducible to the digit.
\emph{(ii) Realistic retrieval}: \texttt{search} returns matched document text, as
production vector stores do; an early variant returning only titles collapsed the
reader to F1 $0.0$ and would have credited the trained store for fixing an
artificially crippled interface. \emph{(iii) Atomic documents}: the store is its
own chunking, so organization is expressible \emph{only} through links and navigation
nodes---any efficiency gain is attributable to structure, and any authored content is
provenance-flagged. \emph{(iv) Contamination control}: we measured gpt-5-mini
answering HotpotQA-style questions at F1 $\approx 1.0$ with \emph{no retrieval at
all}; on public corpora, a reader's score conflates the store with the weights, and
``store value'' loses meaning. A fictional universe makes the no-store score
$\approx 0$, so every point of reader performance is earned through the store.

\section{Training Protocol}
\begin{algorithm}[t]
\caption{Training a knowledge base (reader and editor tool sets as defined
in Sec.~\ref{sec:actions})}
\label{alg:train}
\begin{algorithmic}[1]
\State evaluate reader on \texttt{eval} \Comment{epoch-0 baseline (= B1)}
\For{epoch $= 1 \dots E$}
  \For{$(q, a^{*})$ in shuffled \texttt{train}}
    \State reset short-term memory; task $\gets q$
    \For{step $= 1 \dots N_1$}  \Comment{forward}
      \State $d \gets$ policy(reader tools)
      \If{$d = \texttt{answer}(\hat a)$}
        \State reveal $a^{*}$, F1$(\hat a, a^{*})$; \textbf{break}
      \EndIf
    \EndFor
    \For{step $= 1 \dots N_2$}  \Comment{backward}
      \State $d \gets$ policy(editor tools)
      \If{$d = \texttt{done}$} \textbf{break} \EndIf
    \EndFor
    \State refresh dirty embeddings; log
  \EndFor
  \State snapshot; evaluate reader on \texttt{eval} \Comment{curve point}
\EndFor
\State exam: frozen store, test splits, budget $M$
\end{algorithmic}
\end{algorithm}

Phase 1 uses exactly the reader's toolset and budget, so train-forward accuracy is
directly comparable to test accuracy. The single-shot \texttt{answer} reveals the gold
in its result---standard supervised learning: predict, then see the label, then
update. Leftover Phase-1 budget is forfeited (no smuggling forward steps into
consolidation); a budget-exhausted forward scores $0$ but still receives the gold, so
failed questions---where repair matters most---get targeted consolidation.
Every intermediate store state is reconstructible by replaying the edit trace
(validated byte-exact against epoch snapshots), giving per-iteration store
trajectories for analysis at no storage cost.

\begin{table}[t]
\caption{Credit assignment: Phase-2 diagnosis and repair, determined by the Phase-1
trajectory plus the revealed gold. No oracle localization is provided.}
\label{tab:credit}
\centering
\small
\renewcommand{\arraystretch}{1.15}
\begin{tabular}{L{2.5cm}L{5.2cm}}
\toprule
Trajectory diagnosis & Repair \\
\midrule
Needed fact never surfaced & directed re-search \emph{with the answer in hand}; link
from where the blind search wandered to where the fact lives \\
Fact read but chain unassembled & materialize the chain as links; optionally author a
navigation document for the class \\
Conflicting documents found & resolve using the gold as arbiter (delete or edit the
loser) \\
Class-level gap (same failure would recur on siblings) & build/extend an index document
gathering same-class material \\
\bottomrule
\end{tabular}
\end{table}

The training agent's objectives, stated as goals (mechanisms are its choice):
(1)~answer correctly; (2)~consolidate verified reasoning into structure; (3)~organize
for the question's \emph{class}---build or extend index/navigation documents so unseen
instances of the class resolve in fewer steps; (4)~keep the store
\emph{parsimonious}: duplicated facts compete in search and bury each other; delete
redundancy you create; if the store already serves the class, change nothing.
Objective~(4) encodes the lesson of preliminary runs on an earlier
environment variant, in which an agent trained \emph{without} it quintupled
an already-organized store and regressed the reader by 17 F1 points while
every individual edit looked constructive. Credit assignment needs no oracle: a wrong answer plus the
revealed gold lets the agent re-search \emph{with the answer in hand}, and the repair
is determined by the trajectory difference.

\section{Experimental Design}
\subsection{Benchmarks}
Three arms share one protocol and differ in document granularity and realism
(Table~\ref{tab:arms}).

\begin{table}[t]
\caption{Benchmark arms. All use the same splits scheme, reader, and budgets.}
\label{tab:arms}
\centering
\scriptsize
\renewcommand{\arraystretch}{1.2}
\begin{tabular}{L{1.35cm}L{1.75cm}L{1.75cm}L{1.75cm}}
\toprule
 & \kbgym{} (main) & PhantomWiki (official) & MuSiQue (real text) \\
\midrule
documents & 5{,}864 sentences & 3{,}403 sent.\ or 405 articles & $\sim$11k paragraphs \\
source & our generator & their generator & Wikipedia \\
categories & 10 (26 templates) & 4 shapes & 2--4-hop shapes \\
questions & 330 (150/100/50/30) & 330, same scheme & same scheme \\
supports & exact & unavailable & official \\
contamination & none & none & delta-corrected \\
status & built; main runs & data built & planned \\
\bottomrule
\end{tabular}
\end{table}

\textbf{\kbgym{} (main).} 500 fictional people (10\% name distractors; family
trees, spouses, friendships; job/hobby/city attributes; globally distinct
birthdates) render into 5{,}864 atomic documents. Ten question categories over
26 templates: forward chains QC1--QC3; aggregation counts QC4/QC8 (children,
friends, grandchildren, attribute-filtered counts); abstention QC5;
multi-constraint joins QC6 (no name anchor; unique answer enforced against
distractors); set intersection QC7 (common friends); comparison/superlatives
QC9 (birthdates; ties excluded by construction); reverse lookup QC10 (exactly
one qualifier). Golds are graph-derived with accepted variants; every question
carries its support set. Splits are instance-disjoint: \texttt{train} 150
(i.i.d.\ over all categories, 60\% hard), \texttt{test\_in} 100 (unseen
instances of trained templates), \texttt{test\_out} 50 (one reserved template
per category), \texttt{eval} 30 (stratified; drives the learning curve so the
final test is touched exactly once).

\textbf{Official PhantomWiki (external validity).} A generator we do not
control (phantom-wiki 1.0.3, SWI-Prolog 10.0.2): 16 family trees, 405 person
articles, 480 generated QA over 8 templates grouped into four composition
shapes, with per-question hop depth 1--5; two reserved templates (relational
and aggregative) form \texttt{test\_out}. Their Prolog-derived relations have
no per-fact supports, so repair diagnostics are unavailable on this arm; the
train/test protocol itself does not need them. Documents at sentence or
article granularity (the latter tests the method at paragraph scale with zero
contamination).

\textbf{MuSiQue (real-text arm)~\cite{musique}.} Paragraphs are the documents; the official
composition shapes are the categories, a held-out shape is
\texttt{test\_out}, and the annotated supporting paragraphs restore support
diagnostics. Real corpora are parametrically contaminated (we measured
near-ceiling no-retrieval scores for related datasets), so all scores on this
arm are reported as deltas over a mandatory no-store reader baseline.

\subsection{Baselines}
All baselines are \emph{unsupervised}: structure is induced from the corpus
alone, with no access to questions or answers---the defining contrast with
our trainer, for which gold answers are labels and the structure built is
whatever those labels showed to be missing. Beyond that they differ
\emph{only} in how the store is prepared; reader, test
sets, and budgets are identical, and each reports its one-time build cost
(Table~\ref{tab:baselines}).

\begin{table}[t]
\caption{Baseline arms on \kbgym{}-500. Identical reader, test sets, and
budgets; arms differ only in store preparation.}
\label{tab:baselines}
\centering
\scriptsize
\renewcommand{\arraystretch}{1.2}
\begin{tabular}{L{1.62cm}L{1.5cm}L{2.15cm}L{1.3cm}}
\toprule
Arm & Preparation & Structure added & Build cost \\
\midrule
B1 flat & none & --- & 0 \\
B3 GraphRAG & offline, unsup. & 135 summary docs, 8.5k links & 488k tok \\
B4 RAPTOR & offline, unsup. & summary hierarchy & --- \\
B5 HippoRAG2 & offline, unsup. & 660 index docs, 196k links & $\approx$0 \\
Ours & QA-supervised & links + nav docs & training run \\
\bottomrule
\end{tabular}
\end{table}
\textbf{B1, flat store}: the untrained store (= epoch~0). \textbf{B3,
GraphRAG-style}~\cite{graphrag}: entities extracted per document (lexical
extraction, justified by template-shaped text), co-occurrence communities,
one LLM-written community summary document per community linked to its
members. \textbf{B4, RAPTOR-style}~\cite{raptor} (extended version):
recursive embedding clustering with LLM summary documents per cluster,
hierarchically linked. \textbf{B5, HippoRAG-2-style}~\cite{hipporag2}:
(subject, relation, object) triples per document, entity-sharing links
between documents, per-entity hub documents. B3--B5 structure lands as
authored documents and links \emph{in our store format}, so the common
budgeted reader can traverse it: we test their \emph{structure} under a fixed
reader rather than their retrieval algorithms (a documented adaptation; e.g.\
B5 forgoes Personalized PageRank).

\subsection{Implementation}
Trainer and reader are gpt-5-mini (temperature 0.3, 4{,}096 tokens/turn),
tool-forced decisions, static schemas. Budgets $N_1{=}N_2{=}M{=}15$; FIFO
$K{=}30$; \texttt{search} returns top-5. Embeddings: chroma default ONNX
MiniLM family, environment-maintained. Measured training cost:
$\approx$23k tokens and 83\,s per iteration on the 120-person universe;
the 500-person main run averaged 37.7k tokens and 121\,s per iteration
(11.3M tokens, 10.1\,h, $\approx$\$8 total at current pricing). Per-epoch snapshots plus byte-exact
trace replay reconstruct every intermediate store state for analysis
(merge events replay deterministically without the judge). Basic version:
single seed (seed~0); seed replication planned. The main run reported below has the guard enabled: the judge fired
$\approx$30 times and declined every merge, so at this scale its effect is
null and prompt-level parsimony carries deduplication.

\subsection{Experiment Suite and Predictions}\label{sec:suite}
\textbf{E1 (headline).} Learning curve: eval-split F1 and steps after each
epoch (epoch~0 = B1); final \texttt{test\_in}/\texttt{test\_out} pass on the
last snapshot, at $M \in \{8, 15\}$. Prediction: non-inferior accuracy with
fewer steps, and a leftward shift of the accuracy--budget curve (the trained
store dominates at tight budgets). \textbf{E2.} The baseline table: F1,
steps/question, build tokens, per-question tokens, for B1/B3/B4/B5/ours on
both test splits. \textbf{E3} (log slicing). Three-layer generalization
(train-forward from epoch~2 / \texttt{test\_in} / \texttt{test\_out}) and the
per-category profile. Prediction: forward classes QC1--QC3 are saturated near
1.0 at $M{=}15$ (they are the non-inferiority floor, not the headroom);
gains concentrate in aggregation (QC4/QC8), joins (QC6), intersection (QC7)
and reverse lookup (QC10), whose answers live in relations between documents;
QC5 abstention is expected unchanged. \textbf{E4} (trace mining). (i)
edit-type mix per epoch; (ii) store-size trajectory from replay (the
parsimony objective predicts sublinear growth); (iii) morphology of authored navigation documents
(out-degree distribution, what they gather); (iv) structure--semantics
alignment: the share of built links connecting documents that co-occur in
some support set (supports are invisible to the agent, making the measure
honest); (v) mechanism attribution: correlation between a reader's step
savings and its hits on authored navigation documents and trained links.
Break-even $N^{*}$ (Eq.~\ref{eq:breakeven}) closes the analysis.

\section{Experiment Results}\label{sec:prelim}
All results below are from the final environment (Sec.~\ref{sec:env})
and the main run, keyed to the suite of Sec.~\ref{sec:suite}:
the calibration that fixed the main configuration
(Sec.~\ref{sec:res-calib}), then the learning curve (E1,
Sec.~\ref{sec:res-e1}), the baseline table (E2, Sec.~\ref{sec:res-e2}),
retention and its per-category profile (E3, Sec.~\ref{sec:res-e3}), and
trace mining (E4, Sec.~\ref{sec:res-e4}).

\begin{figure}[t]
\centering
\includegraphics[width=\columnwidth]{learning.pdf}
\caption{Main run learning signals. (a)~Held-out eval by epoch (30 questions
per point: 20 trained-template, 10 reserved-template). (b)~Rolling
(window~30, drawn from the first full window) forward accuracy over
training; the dashed line marks epoch~2, where questions are re-encountered
and retention is carried only by the store.}
\label{fig:learning}
\end{figure}

\begin{figure*}[t]
\centering
\includegraphics[width=0.92\textwidth]{network.pdf}
\caption{Evolution of the document network during training, every 60
iterations. Every document is drawn; position is a t-SNE projection of its
embedding, so proximity = semantic similarity (the clusters visible at
iteration~0 are the universe's natural topic groups). Orange: agent-authored
navigation documents (size = link degree); green: links. Structure grows
from nothing into a navigation layer spanning the semantic space: hub
directories (dense orange cluster) fan out into topic clusters, and local
links connect semantically adjacent documents.}
\label{fig:network}
\end{figure*}

\begin{figure}[t]
\centering
\includegraphics[width=0.8\columnwidth]{trajectory.pdf}
\caption{Store trajectory during training, reconstructed by trace replay.
\emph{Authored documents} = agent-written documents still alive and unedited
(navigation/index documents); \emph{net documents added} = total store size
minus the initial 5{,}864 (authored survivors, including later-edited ones,
minus deleted originals). Growth is steady and small ($+6.6\%$ over 300
iterations).}
\label{fig:trajectory}
\end{figure}

\subsection{Calibration and Cost Profile}\label{sec:res-calib}
Reader difficulty was calibrated by probes on untrained stores: forward named-entity
chains are natively easy for document-level retrieval at any tested scale (F1
$1.0$); headroom concentrates in the relation-level classes and grows with universe
size (untrained F1 $0.83$ at 120 people vs.\ $0.67$ at 500, mean steps $6.7$ vs.\
$9.2$)---motivating the 500-person main configuration, where the cost axis has
roughly $3\times$ headroom over the $\sim$3-step floor. Accuracy can
saturate while cost headroom persists---one reason efficiency is the
primary endpoint. The 300-iteration
main run consumed 11.3M training tokens over 10.1\,h ($\approx\$8$ at
current pricing).

\subsection{E1: Training Dynamics and the Learning Curve}\label{sec:res-e1}
The main run (500-person universe, full protocol; 300 iterations, 11.3M
training tokens, 10.1h) yields the training half of E1; the frozen-store
exam that completes it is reported below at $M{=}15$ and $M{=}8$. Two
training-time signals, then the exam.

\textbf{Held-out eval drifts up.} Eval-in $0.60 \to 0.60 \to 0.65$,
eval-out $0.50 \to 0.60 \to 0.60$ across epochs 0/1/2
(Fig.~\ref{fig:learning}a; $n{=}30$ per point, so these are directional).

\textbf{The store improves \emph{within} an epoch.} Rolling forward
accuracy climbs through epoch~1 (first 50 iterations $0.56$, last 50
$0.64$): later training questions already benefit from structure
consolidated by earlier ones, before any question repeats. This is the
evidence that the structure is genuinely accumulating rather than being
rebuilt per question: what earlier questions left behind is already paying
off before anything repeats.

\textbf{The frozen exam: accuracy transfers; efficiency does not.} An
unchanged reader on the frozen epoch-2 store is non-inferior to B1 at both
budgets, and at neither budget is it cheaper. At $M{=}15$ it scores
$0.630$/$0.660$ (\texttt{test\_in}/\texttt{test\_out}) vs.\ B1's
$0.620$/$0.620$ ($\Delta$ pooled $+0.020$; all $\Delta$ in this
paper are trained minus B1, the untrained flat store, on the same
questions), taking $8.73$ vs.\ $8.97$ steps
($\rho = 0.97$) and saving $80$ tokens per question. At $M{=}8$ it scores
$0.570$/$0.540$ vs.\ B1's $0.530$/$0.600$ ($\Delta$ pooled $+0.007$: better
on trained templates, \emph{worse} on reserved ones), at $6.23$ vs.\ $6.21$
steps ($\rho = 1.00$) and $194$ tokens \emph{more}. Both budgets clear the
non-inferiority margin ($-\delta = -0.03$ pooled), so the store the agent
built is not damaged---but the efficiency hypothesis, our central claim, is
\textbf{not supported on unseen questions}: $\rho \approx 1$ at both
budgets, break-even $N^{*} \approx 1.4 \times 10^{5}$ questions at $M{=}15$
and undefined at $M{=}8$, and the predicted leftward shift of the
accuracy--budget curve does not appear (the trained and flat curves are
within noise of each other at both points). The $-38\%$ step saving of
Sec.~\ref{sec:res-e3} is real but confined to \emph{re-encountered}
questions. Sections~\ref{sec:res-e2} and~\ref{sec:res-e4} locate the gap:
no baseline saves steps either, and the structure that was built pays only
where derivation is expensive.

\begin{table}[t]
\caption{E2: offline-construction baselines under the common reader
(final environment, gpt-5-mini). F1 and steps per question on
\texttt{test\_in} ($n{=}100$) and \texttt{test\_out} ($n{=}50$); tok/q =
mean answering tokens across both splits; Build = one-time construction
cost. Best per column in bold.}
\label{tab:e2}
\centering
\scriptsize
\renewcommand{\arraystretch}{1.15}
\setlength{\tabcolsep}{3.5pt}
\begin{tabular}{llrrrrr}
\toprule
 & & \multicolumn{2}{c}{\texttt{test\_in}} & \multicolumn{2}{c}{\texttt{test\_out}} & \\
\cmidrule(lr){3-4}\cmidrule(lr){5-6}
Store & Build & F1 & steps & F1 & steps & tok/q \\
\midrule
\multicolumn{7}{l}{\emph{budget $M{=}15$}}\\
B1 flat & --- & 0.620 & 9.0 & 0.620 & 8.9 & 9.7k \\
B3 GraphRAG-style & 488k tok & 0.570 & 9.3 & 0.640 & 9.1 & 10.5k \\
B5 HippoRAG2-style & $\sim$0 & \textbf{0.660} & 9.3 & \textbf{0.700} & 7.9 & 11.1k \\
Ours (trained) & 11.3M tok & 0.630 & 9.0 & 0.660 & 8.3 & 9.7k \\
\midrule
\multicolumn{7}{l}{\emph{budget $M{=}8$}}\\
B1 flat & --- & 0.530 & 6.5 & \textbf{0.600} & 5.7 & 4.9k \\
B3 GraphRAG-style & 488k tok & 0.560 & 6.2 & 0.500 & 6.2 & 5.0k \\
B5 HippoRAG2-style & $\sim$0 & 0.530 & 6.4 & 0.520 & 5.9 & 6.0k \\
Ours (trained) & 11.3M tok & \textbf{0.570} & 6.3 & 0.540 & 6.1 & 5.1k \\
\bottomrule
\end{tabular}
\end{table}

\subsection{E2: Baselines Buy Accuracy, Not Steps}\label{sec:res-e2}
Table~\ref{tab:e2} reports the offline-construction baselines under the
common reader. Three observations. \textbf{(i) No baseline saves steps.}
Every arm sits at $9.0$--$9.3$ steps/question at $M{=}15$: community
summaries and entity indexes change \emph{what} the reader finds, not how
long it searches, so cost reduction---our primary endpoint---is not a free
by-product of offline structure. \textbf{(ii) Entity indexing is the
accuracy frontier.} B5 leads both splits ($0.660$/$0.700$) at essentially
zero build cost, making it the accuracy bar for the trained store.
\textbf{(iii) Summaries can pollute.} B3 spends a 488k-token build to
\emph{lose} $0.05$ F1 on \texttt{test\_in} relative to the flat store: its
135 community summaries compete with source documents in retrieval and bury
them---summary pollution arising inside a standard pipeline. At the tight budget ($M{=}8$) all stores degrade and the ordering
compresses; no arm dominates. The trained store enters the table without
dominating it: at $M{=}15$ it is above B1 and B3 on both splits and below
B5's entity indexes; at $M{=}8$ it leads \texttt{test\_in} but trails B1 on
\texttt{test\_out}. Crucially it shares the baselines' defining failure ---
$\rho \approx 1$ against B1 at both budgets. That \emph{every} method here,
supervised or not, lands within noise of the flat store
on steps is the paper's most robust negative finding, and the step
distribution says why.

\textbf{Retrieval cannot enumerate, and that is what structure is for.}
The reader's \texttt{search} returns the five most similar documents and
nothing else: no paging, no count, no way to ask for the sixth. A question
whose answer requires a \emph{set} --- how many grandchildren, which friend
is shared, which holder of a job lives in a city --- is therefore not
merely hard for such a reader, it is out of reach unless the store itself
makes the set reachable, because the twenty-five residents of a city cannot
be recovered five at a time from a store that does not connect them. This
predicts exactly where offline structure should pay, and the baselines
confirm it. Partitioning our ten classes into those that need a set
(QC4, QC6--QC8, QC10) and those answerable from a single document or a
chain (QC1--QC3, QC5, QC9), B5's entity indexes lift the first group from
$0.388$ to $0.514$ F1 and the second group from $0.797$ to $0.797$ --- a
gain of $+0.126$ where enumeration is required and of exactly zero
where it is not. Per class the pattern is sharp: counts QC4 $+0.30$,
intersection QC7 $+0.29$, join QC6 $+0.09$, while QC1 and QC2 sit at
$1.00$ under every arm. Structure in a store is not a general accelerant;
it is specifically a substitute for the enumeration that vector retrieval
cannot perform.

\textbf{Why no method saves steps: mean steps is not a lookup-length
measure.} Steps per question are strongly bimodal --- a mass around 3--6
and a spike at exactly $M$ --- because a question the reader cannot resolve
consumes the \emph{entire} budget by construction. The mean therefore
decomposes as $\bar{s} = p_{\text{give-up}} M + (1 - p_{\text{give-up}})
\bar{s}_{\text{answered}}$, and at $M{=}15$ the second term is nearly
arm-invariant: $6.77$ (B1), $6.97$ (B3), $6.91$ (B5), $6.90$ (ours). The
cleanest contrast is the paired one: on the 103 questions that
\emph{both} B1 and the trained store answered, the trained store takes
$6.48$ steps against B1's $6.57$ --- a $0.10$-step difference, i.e.\ none.
What the trained store actually does is convert give-ups: it rescues 13
questions B1 abandoned (answering them in $10.2$ steps at F1 $0.69$) while
losing 7 in the other direction, and that net of 6 fully accounts for the
mean dropping from $8.97$ to $8.73$. Organization thus moves the
\emph{give-up rate}, not the length of a successful lookup. Two
consequences. Methodologically, mean steps conflates two different
quantities and is a poor efficiency endpoint at budgets where give-ups are
common ($27$--$35\%$ here); reporting $p_{\text{give-up}}$ and
$\bar{s}_{\text{answered}}$ separately, as we do above, is the honest
decomposition. Substantively, the reader's stopping rule is
confidence-driven rather than structure-driven: it stops when it believes
it can answer, so a shortcut it traverses earlier does not make it stop
earlier --- it only changes \emph{whether} it ever gets there. Making
materialized structure pay in steps therefore requires a reader that can
recognize it is on a shortcut, which our fixed reader, chosen deliberately
to isolate the store, cannot.

\subsection{E3: Retention and Its Per-Category Profile}\label{sec:res-e3}
\textbf{Retention is large, and concentrated exactly where predicted.}
On re-encountered questions (epoch~2) forward F1 reaches $0.85$
(vs.\ $0.59$ in epoch~1). Table~\ref{tab:retention} decomposes by category:
forward chains were saturated from the start (QC1--QC3, QC9 at
$0.86$--$1.00$ in both epochs --- retrieval alone suffices), while the
relation-level classes, whose answers live \emph{between} documents, jump:
joins QC6 $0.46 \to 1.00$, intersection QC7 $0.62 \to 1.00$, reverse lookup
QC10 $0.25 \to 0.83$, abstention QC5 $0.00 \to 0.58$, counts QC4
$0.17 \to 0.50$ and QC8 $0.46 \to 0.58$. This is the design prediction of E3
(Sec.~\ref{sec:suite}) realized: training gains concentrate in the classes retrieval
cannot solve natively.

\textbf{Retention is also cheaper, with a caveat about what is being
compared.} On re-encountered questions the forward pass takes $5.7$ steps
and 31k tokens, vs.\ $9.2$ steps and 45k tokens in epoch~1: $-38\%$ steps,
$-32\%$ tokens. Decomposing as in Sec.~\ref{sec:res-e2}, both terms move:
budget exhaustion falls from $47/150$ ($31\%$) to $7/150$ ($5\%$), and on
the 100 questions resolved within budget in \emph{both} epochs the lookup
itself shortens from $6.43$ to $4.64$ steps ($-28\%$) while F1 rises from
$0.870$ to $0.950$. Unlike the null result on unseen questions, then, the
conversion of derivation into lookup is genuine here. The caveat is that
these two epochs are not a controlled store contrast: the store evolves
\emph{during} epoch~1, so an early epoch-1 question meets a nearly
untrained store and a late one meets a partly trained store, and the epoch
label mixes store state with question order. The controlled version --- the
same 150 training questions examined under the frozen initial store and the
frozen final store, exactly as the test splits are examined --- is reported
in Sec.~\ref{sec:res-gap}.

\textbf{Three-layer decomposition, completed.} Retention $0.85$ /
\texttt{test\_in} $0.63$ / \texttt{test\_out} $0.66$. Both test layers
sit above B1, so the trained organization transfers accuracy --- but the
$0.2$-point gap to retention says most of the built structure is
instance-bound rather than class-general. The per-category transfer
profile ($\Delta$ vs.\ B1 at $M{=}15$, both splits pooled, same questions) is
correspondingly modest: QC4 $+0.10$, QC9 $+0.08$, QC3/QC7 $+0.07$, QC6
$+0.04$, QC5 $0.00$, QC8 $-0.05$, QC10 $-0.12$ ($n{=}8$). The sharpest
contrast is reverse lookup: the class with the largest retention jump
($0.25 \to 0.83$) transfers \emph{worst} --- its navigation documents
memoize trained instances, not the class.

\begin{table}[t]
\caption{Per-category forward F1: first encounter (epoch 1) vs.\
re-encounter (epoch 2). Retention concentrates in relation-level classes.}
\label{tab:retention}
\centering
\scriptsize
\renewcommand{\arraystretch}{1.15}
\begin{tabular}{L{2.3cm}rrr}
\toprule
Category & $n$ & epoch 1 & epoch 2 \\
\midrule
QC1--QC3, QC9 (chains, compare) & 55 & 0.86--1.00 & 0.86--1.00 \\
QC6 join & 26 & 0.46 & \textbf{1.00} \\
QC7 intersection & 13 & 0.62 & \textbf{1.00} \\
QC10 reverse lookup & 12 & 0.25 & \textbf{0.83} \\
QC5 abstention & 12 & 0.00 & \textbf{0.58} \\
QC4 counts & 6 & 0.17 & \textbf{0.50} \\
QC8 deep counts & 26 & 0.46 & 0.58 \\
\bottomrule
\end{tabular}
\end{table}

\subsection{The Controlled Generalization Gap}\label{sec:res-gap}
The epoch-to-epoch contrast above mixes store state with question order. The
controlled comparison holds the store frozen and varies only which questions
are asked: the same reader, budget and protocol as the test splits, applied
to the 150 \emph{training} questions under (a) the untrained initial store
and (b) the frozen epoch-2 store. Together with the test-split results this
gives the train/test generalization gap in its standard form --- performance
on questions the store was optimized for, against performance on questions
it was not, with everything else held fixed. This exam was running at
submission time; the slot is reserved here because it is the measurement
that separates an overfitted store from a generalizing one, and the
evidence so far (large retention gains, near-null transfer of the step
saving) predicts a wide gap.

\subsection{E4: Trace Mining---Edit Behavior and Store Morphology}\label{sec:res-e4}
\textbf{Where the structure lives.} Figure~\ref{fig:network} maps the
store's evolution across its semantic space (every document drawn; position
= t-SNE of the document embedding, so the clusters visible at iteration~0
are the universe's natural topic groups). Navigation documents appear
throughout the map; hub directories --- the dense orange cluster --- fan out
into the topic clusters, and short local links connect semantically adjacent
documents: the store visibly acquires an access layer it was born without.
Figure~\ref{fig:trajectory} quantifies the same growth from trace replay:
$+390$ documents ($+6.6\%$), 597 links, accumulating steadily across both
epochs --- the sublinear, structural growth the parsimony objective was
designed to enforce.

\textbf{Edit behavior.} The backward pass writes $\approx$1.6 documents and
$\approx$2.0 links per iteration in epoch~1; in epoch~2 the add rate falls
to $1.0$/iteration while linking stays at $2.0$ --- consistent with the
parsimony objective (less to author once the structure exists). Edits and
deletes stay rare ($0.13$ and $0.03$/iteration). The dedup judge fired
$\approx$30 times and correctly declined every candidate (zero merges):
prompt-level parsimony already suppressed duplication at this scale.

\textbf{Structure--semantics alignment: the links are real.} Of the 596
built links, 469 attach a navigation document (no support-set analogue
exists for those) and 127 connect two statements. Of those 127, $25.2\%$
join documents that co-occur in the support set of some universe
question---against a $0.03\%$ chance rate for random statement pairs, a
$720\times$ lift. Support sets are never visible to the training agent, so
this is an honest measure: when the agent links two facts, it is recovering
genuine reasoning adjacency, not surface similarity.

\textbf{Mechanism attribution: the observational split does not identify,
and we say so.} We had intended to attribute the reader's step savings to
its hits on built structure, splitting exam questions by whether the
trajectory touched an authored navigation document or a linked statement.
That split is not identifying, for a mechanical reason: touching is
\emph{downstream} of searching. The probability that a trajectory
encounters built structure rises from $0\%$ when the reader issues one
search to $63\%$ when it issues nine, so the touched group is selected for
long trajectories by construction. Difficulty matching does not repair
this: the touched group remains $+0.92$ steps worse than the untouched one
after exact stratification on B1's step count, which is exactly what
reverse causation predicts. Two secondary measurements rule out the
explanation we first reached for --- retrieval crowding. Authored documents
occupy $5.3\%$ of top-5 retrieval slots against a $5.7\%$ share of the
store by count, i.e.\ no over-retrieval, and their share \emph{rises} with
question difficulty ($2.2\%$ easy, $6.6\%$ hard), which is what relevance,
not pollution, looks like.

\textbf{What the agent actually wrote.} Inspection of the 359 authored
documents separates two kinds. Most ($80\%$) materialize an answer:
\emph{``Caleb Netherfield's father is Caleb Nightingale and Caleb
Nightingale's spouse is Petra Jessup and Petra Jessup's job is joiner''}
collapses a three-hop chain into one document; others aggregate a reverse
index (\emph{``People whose hobby is pottery include \ldots''}) or a count.
The remaining $20\%$ are pointer documents that state no fact at all ---
\emph{``Spouse job index gathers notes that record the jobs of spouses of
people''}, \emph{``A Junia Wetherby directory links to Junia Wetherby's
hobby, job, residence, friends, parent, and birth date notes''}. A reader
that retrieves one of these pays a slot and a read step and learns nothing,
then still has to find the facts. Notably the linking effort went to the
wrong kind: pointer documents carry $2.11$ links on average against $1.06$ for
contentful ones. Their absolute retrieval footprint is small ($0.9\%$ of
slots), so this is a design smell rather than a proven cost --- which is
precisely why the causal question needs an intervention on the store rather
than an observational split.

\textbf{The store never became navigable.} The action set can produce
exactly two kinds of structure: index documents that point at the documents
belonging to their key, and edges between documents that already exist
(fact-to-fact, or index-to-index, which is how navigation acquires
levels). Profiling the trained store by those classes explains the null
transfer better than any analysis of the reader's behaviour. Of the 359
authored index documents, the mean out-degree is $1.27$ and the median is
$1$; \textbf{151 of them ($42\%$) carry no links at all}, and 340 ($95\%$)
carry fewer than five. The store holds 316 index-to-fact edges, 139
index-to-index edges and 127 fact-to-fact edges. The consequence is
stark: \textbf{only 247 of 5{,}891 facts ($4.2\%$) are reachable from any
index in a single read}. A store in which $96\%$ of the content has no
entry point is not an organized store carrying instance-bound structure ---
it is an unorganized store with 359 mostly empty labels attached. The
efficiency hypothesis was therefore never really tested by this run: the
mechanism it posits requires navigable structure, and the agent did not
build any. This relocates the finding from ``materialized structure does
not transfer'' to ``this training procedure does not produce materialized
structure'', which is a statement about the optimizer, not about the
hypothesis --- and a fixable one. Sec.~\ref{sec:disc} takes it up.

\textbf{The ablation that would settle it.} Holding reader, questions and
budget fixed and removing structure from the store gives the missing
counterfactual: the trained store minus its 359 authored documents (127
statement-to-statement links retained), and minus its links (authored
documents retained). Both exams were running at submission time. This is
the design the confound above forces, and we report the confound rather
than the attribution it invalidates.

\emph{Caveats: single seed; per-class $n$ as low as 6; epoch-2 ordering is
reshuffled, so per-class retention rides on small samples. The frozen-store exam separates transferable organization from instance memoization.}

\section{Discussion and Limitations}\label{sec:disc}
\textbf{Curation is a boundary, not a slogan.} Our results suggest
curation creates value where organization is lacking and pollution where it
is not: training helps on the disorganized main condition, while
GraphRAG-style summaries added to an already-searchable store regress
\texttt{test\_in} (Sec.~\ref{sec:res-e2}). The difference is carried by
whether the optimizer knows when \emph{not} to edit.
The train/test protocol makes this boundary measurable; downstream-only evaluation
cannot see a store being slowly ruined while individual answers still look fine.

\textbf{Reading the negative result.} \emph{The proximate cause is
construction, not concept.} Sec.~\ref{sec:res-e4} shows the trained store
never became navigable --- $42\%$ of its index documents point at nothing and
$96\%$ of its facts have no entry point --- so the run tested a procedure
that failed to build indexes rather than the value of indexes. Three
plausible reasons, each actionable. The objectives named the goal
(``organize for the question's class'') without naming the mechanism, so
the cheapest compliant act was to record the verified answer, which is
what $80\%$ of authored documents do. The backward budget (15 actions) cannot
populate a thirty-member index in one iteration, and nothing in the
protocol asked the agent to \emph{return} to an index and extend it, so
indexes were created and abandoned. And the parsimony objective, added to
stop the pollution of Sec.~\ref{sec:design}, penalizes adding documents
without distinguishing a redundant fact from an entry point that is
supposed to accumulate links. A revised protocol that names indexes as the
unit of construction, requires extension of an existing index over
creation of a new one, and scores an index by its link completeness is the
obvious next experiment; we state it as the concrete next step rather than
claim the hypothesis is settled.

\textbf{The broader reading.} Our central hypothesis---that
materialized structure lets an unchanged reader answer unseen questions in
fewer steps---is not supported at the budgets we tested. Three observations
keep this informative rather than merely disappointing. First, the effect
\emph{does} exist within the training distribution: re-encountered
questions run $38\%$ shorter, so structure genuinely converts derivation
into lookup; what fails is transfer of that saving to new instances.
Second, the failure is not ours alone---GraphRAG- and HippoRAG2-style
offline construction save no steps either, which suggests the bottleneck is
the agentic reader's search policy rather than any particular way of
organizing a store: a reader that decides when to stop searching by its own
uncertainty will not stop earlier just because a shortcut exists. Third,
we cannot yet say which part of what the agent built carries the retained
gain: the observational attribution turned out to be confounded by reverse
causation (Sec.~\ref{sec:res-e4}), and the store ablations that would
answer it are running. Together these point at a concrete next design: gate
index construction --- and index \emph{use} --- on the expected derivation
cost of a question class, rather than indexing whatever was just answered.

\textbf{Supervised, not unsupervised.} Every offline baseline induces its
structure from the corpus alone and can never incorporate what the questions
reveal; our trainer is told, question by question, which answers it failed to
support, and builds against that signal. In machine-learning terms the
baselines are unsupervised pretraining of structure and ours is supervised
learning of it, which is why the gains concentrate in the question classes
the labels actually covered.

\textbf{Deployment reading.} The abstraction targets a real configuration: agent
fleets maintaining shared repositories, wikis, and knowledge bases. Our failure
ledger maps directly onto production incidents---duplicated near-copies burying each
other is the lost-update problem; authored-without-provenance content is the
hallucinated fix; a deleted last instance is knowledge loss during refactoring. The
break-even point $N^{*}$ answers the deployment question in units operators use:
after how many queries does curation pay for itself. Our preliminary runs
(Sec.~\ref{sec:design}) double as a default-behavior audit: an agent given
a maintenance objective but no parsimony objective degraded the asset it was
tending while every individual action looked constructive---a caution for
unsupervised agent operation on shared state.

\textbf{Threats to validity.} Template-rendered language is simpler than natural
prose; phrasing regularity may flatter lexical matching (mitigated by embedding
retrieval and by the planned official-PhantomWiki arm, whose generator we do not
control). The reader and trainer share a model family, so structure tuned by the
trainer may suit that family's search habits. Single seed in the basic version;
single-agent by design (concurrent maintenance, with its conflict taxonomy, is the
natural sequel); support-set diagnostics are unavailable on the external arm.
Grading is token-F1 against short golds---adequate for entity answers, blunt for
free-form ones; the environment avoids the latter by construction.

\textbf{An online variant.} We deliberately separate a training phase from a
frozen-store examination, because that separation is what makes the
generalization question answerable at all: a store edited during evaluation
cannot be said to have been tested on anything. Nothing in the method
requires it. The training loop already answers a question before it edits,
so the same agent could curate against live traffic, treating whatever
feedback a deployment provides---a correction, a follow-up, a click---as the
label the gold answer supplies here. What such a variant would need is not a
different curator but a different measurement: with the store changing
underneath, held-out accuracy stops being well defined, and the honest
alternative is a prequential one, scoring each question before the edits it
provokes.

\begin{thebibliography}{00}
\bibitem{llmwiki-rar} Retrieval as reasoning: self-evolving agent-native retrieval
via LLM-wiki, arXiv:2605.25480, 2026.
\bibitem{llmwiki-stream} Streaming knowledge compilation: proactive
materiality-scored pinning for time-evolving LLM wikis, arXiv:2606.09877, 2026.
\bibitem{karpathy-llmwiki} A.~Karpathy, ``llm-wiki,'' public note, Apr.\ 2026.
\bibitem{graphrag} D.~Edge \emph{et al.}, ``From local to global: a GraphRAG approach
to query-focused summarization,'' arXiv:2404.16130, 2024.
\bibitem{raptor} P.~Sarthi \emph{et al.}, ``RAPTOR: recursive abstractive processing
for tree-organized retrieval,'' ICLR, 2024.
\bibitem{hipporag2} B.~Guti\'errez \emph{et al.}, ``From RAG to memory:
non-parametric continual learning for LLMs,'' ICML, 2025.
\bibitem{ace} ``Agentic context engineering: evolving contexts for self-improving
language models,'' ICLR, 2026.
\bibitem{phantomwiki} A.~Gong \emph{et al.}, ``PhantomWiki: on-demand datasets for
reasoning and retrieval evaluation,'' ICML, 2025.
\bibitem{ircot} H.~Trivedi \emph{et al.}, ``Interleaving retrieval with
chain-of-thought reasoning for knowledge-intensive multi-step questions,'' ACL, 2023.
\bibitem{selfrag} A.~Asai \emph{et al.}, ``Self-RAG: learning to retrieve, generate,
and critique through self-reflection,'' ICLR, 2024.
\bibitem{agentic-rag-survey} ``Reasoning RAG via system 1 or system 2: a survey on
reasoning agentic retrieval-augmented generation,'' arXiv:2506.10408, 2025.
\bibitem{amem} W.~Xu \emph{et al.}, ``A-MEM: agentic memory for LLM agents,''
arXiv:2502.12110, 2025.
\bibitem{longmemeval} D.~Wu \emph{et al.}, ``LongMemEval: benchmarking chat
assistants on long-term interactive memory,'' ICLR, 2025; and LongMemEval-V2,
arXiv:2605.12493, 2026.
\bibitem{memdelta} ``MemDelta: controlled baselines and hidden confounds in agent
memory evaluation,'' arXiv:2606.29914, 2026.
\bibitem{letta-sleep} Letta, ``Sleep-time compute,'' 2025.
\bibitem{rag-lewis} P.~Lewis \emph{et al.}, ``Retrieval-augmented generation
for knowledge-intensive NLP tasks,'' NeurIPS, 2020.
\bibitem{dpr} V.~Karpukhin \emph{et al.}, ``Dense passage retrieval for
open-domain question answering,'' EMNLP, 2020.
\bibitem{realm} K.~Guu \emph{et al.}, ``REALM: Retrieval-augmented language
model pre-training,'' ICML, 2020.
\bibitem{fid} G.~Izacard and E.~Grave, ``Leveraging passage retrieval with
generative models for open domain question answering,'' EACL, 2021.
\bibitem{react} S.~Yao \emph{et al.}, ``ReAct: Synergizing reasoning and
acting in language models,'' ICLR, 2023.
\bibitem{selfask} O.~Press \emph{et al.}, ``Measuring and narrowing the
compositionality gap in language models,'' Findings of EMNLP, 2023.
\bibitem{flare} Z.~Jiang \emph{et al.}, ``Active retrieval augmented
generation,'' EMNLP, 2023.
\bibitem{deeprag} X.~Guan \emph{et al.}, ``DeepRAG: Thinking to retrieval
step by step for large language models,'' arXiv:2502.01142, 2025.
\bibitem{searchr1} B.~Jin \emph{et al.}, ``Search-R1: Training LLMs to reason
and leverage search engines with reinforcement learning,'' arXiv:2503.09516,
2025.
\bibitem{memgpt} C.~Packer \emph{et al.}, ``MemGPT: Towards LLMs as operating
systems,'' arXiv:2310.08560, 2023.
\bibitem{mem0} P.~Chhikara \emph{et al.}, ``Mem0: Building production-ready
AI agents with scalable long-term memory,'' arXiv:2504.19413, 2025.
\bibitem{generative-agents} J.~S. Park \emph{et al.}, ``Generative agents:
Interactive simulacra of human behavior,'' UIST, 2023.
\bibitem{voyager} G.~Wang \emph{et al.}, ``Voyager: An open-ended embodied
agent with large language models,'' arXiv:2305.16291, 2023.
\bibitem{reflexion} N.~Shinn \emph{et al.}, ``Reflexion: Language agents with
verbal reinforcement learning,'' NeurIPS, 2023.
\bibitem{expel} A.~Zhao \emph{et al.}, ``ExpeL: LLM agents are experiential
learners,'' AAAI, 2024.
\bibitem{awm} Z.~Z. Wang \emph{et al.}, ``Agent workflow memory,''
arXiv:2409.07429, 2024.
\bibitem{cheatsheet} M.~Suzgun \emph{et al.}, ``Dynamic cheatsheet: Test-time
learning with adaptive memory,'' arXiv:2504.07952, 2025.
\bibitem{storm} Y.~Shao \emph{et al.}, ``Assisting in writing Wikipedia-like
articles from scratch with large language models,'' NAACL, 2024.
\bibitem{lightrag} Z.~Guo \emph{et al.}, ``LightRAG: Simple and fast
retrieval-augmented generation,'' arXiv:2410.05779, 2024.
\bibitem{hipporag} B.~Guti\'errez \emph{et al.}, ``HippoRAG:
Neurobiologically inspired long-term memory for large language models,''
NeurIPS, 2024.
\bibitem{hotpotqa} Z.~Yang \emph{et al.}, ``HotpotQA: A dataset for diverse,
explainable multi-hop question answering,'' EMNLP, 2018.
\bibitem{musique} H.~Trivedi \emph{et al.}, ``MuSiQue: Multihop questions via
single-hop question composition,'' TACL, 2022.
\bibitem{twowiki} X.~Ho \emph{et al.}, ``Constructing a multi-hop QA dataset
for comprehensive evaluation of reasoning steps,'' COLING, 2020.
\bibitem{squad} P.~Rajpurkar \emph{et al.}, ``SQuAD: 100,000+ questions for
machine comprehension of text,'' EMNLP, 2016.
\bibitem{locomo} A.~Maharana \emph{et al.}, ``Evaluating very long-term
conversational memory of LLM agents,'' ACL, 2024.
\bibitem{rome} K.~Meng \emph{et al.}, ``Locating and editing factual
associations in GPT,'' NeurIPS, 2022.
\bibitem{memit} K.~Meng \emph{et al.}, ``Mass-editing memory in a
transformer,'' ICLR, 2023.
\bibitem{star} E.~Zelikman \emph{et al.}, ``STaR: Bootstrapping reasoning
with reasoning,'' NeurIPS, 2022.
\end{thebibliography}

\end{document}
